%% appendix-complexity -- proofs for the part that is not mechanized
%% (sec:cx): the smaller unit reader and the antecedent rejection clock
%% (sec:un-smaller), leaf cones and tree classes, and signed satisfiability
%% with copy literals (sec:cx-facts).

\section{Proofs for Complexity and Sharper Bounds (Not Mechanized)}\label{app:cx}

This appendix proves the results of \cref{sec:cx} whose proofs were omitted
there. Like that section, it is not mechanized, and nothing outside
\cref{sec:cx} uses it.

\subsection{The Smaller Unit Reader and the Rejection Clock}\label{app:cx-smaller}

\begin{proof}[Proof of \cref{lem:un-reduced}]
(1) The bit $b$ never returns from $1$ to $0$, and if $\mathrm{du}$ fails at a
cut, it fails at every later cut (\cref{lem:un-source}(2)). So the replaced
states are closed under the transitions of $\RU(p)$, and a run that meets one
of them stays among them. None of them is rejected, since a rejected state has
$b=0$ and satisfies $\mathrm{du}$. Hence a word rejected by $\RU(p)$ meets no
replaced state, and a word rejected by $\widehat{\mathcal R}_p$ never meets the
absorbing state $\bot$. In both cases the run consists of kept states, between
which the transitions are unchanged, so the word has the same run in both
readers and is rejected by both. Each comparison of the word is tested on
states of this run, which carry the same self guard in both readers. So the two
readers reject the same words, with the same states at every cut and the same
admitted comparisons, and they have the same holes.

(2) If two prefixes of a word are rejected by $\widehat{\mathcal R}_p$, they
are rejected by $\RU(p)$ by (1); by betweenness for $\RU(p)$
(\cref{lem:un-rejected-path}), so is every intermediate prefix, and by (1)
again it is rejected by $\widehat{\mathcal R}_p$. A word rejected by
$\widehat{\mathcal R}_p$ has, by (1), the same first rejected cut and the same
admitted comparisons in both readers, so the admission cone of $\RU(p)$
transfers.

(3) A kept state after a word $u$ is
$\zeta=(P^\varphi_p(u),0,P^\psi_p(u))$ by \cref{lem:un-source}(1), and
$\mathrm{du}(\zeta)$ gives $P^\psi_p(u)\cap\DeadU^\eps_\psi\ne\varnothing$, so
$P^\psi_p(u)\ne\varnothing$. By \cref{lem:sm-duality}(3), applied in $\varphi$
and in $\psi$ to the same seed $p$ and the same word $u$, all elements of
$P^\varphi_p(u)$ and of $P^\psi_p(u)$ have the sign $\sgn(p)\oplus\parity(u)$,
which the nonempty set $P^\psi_p(u)$ determines. For each sign there are at
most $2^v$ sets $P\subseteq X^\pm$ and $2^s-1$ nonempty sets
$P'\subseteq(X\cup Z)^\pm$ of that sign. Adding $\bot$ gives at most
$1+2^{v+1}(2^s-1)=2^{v+s+1}-2^{v+1}+1<2^{v+s+1}$ states.
\end{proof}

\begin{proof}[Proof of \cref{lem:un-clock}]
Let $w$ be rejected by $\RU(p)$, and let $\zeta_j=(P_j,b_j,P'_j)$ be the state
of $\Bsrc_p$ at the cut $j$, so that $P_j=P^\varphi_{p,j}(w)$ is also the state
of $\mathcal A_p$ there. At the last cut, $b=0$ and $\mathrm{du}$ holds, so by
\cref{lem:un-source}(2) and the monotonicity of the bit, $b_j=0$ and
$\mathrm{du}(\zeta_j)$ hold at every cut $j$. Hence $\zeta_j\in\Rej_{\mathrm U}$
iff $\mathrm{dt}(\zeta_j)$ fails, iff
$P_j\cap\DeadT^\eps_\varphi=\varnothing$, which is rejection by $\mathcal A_p$
at the cut $j$. So $w$ is rejected by $\mathcal A_p$, with the same first
rejected cut, and $r(w)$ is the same for both readers. By sign coherence
(\cref{lem:sm-duality}(3)), every reachable state of $\mathcal A_p$ is empty or
a nonempty set of elements of $X^\pm$ of one sign, so $\mathcal A_p$ has at
most $2(2^v-1)+1=2^{v+1}-1$ states. \Cref{thm:bg-depth} holds for every reader
and every nonempty rejected word, without any condition on the admission
relation, so for $w\ne\eps$ it gives
$\depth(w)\le2(2^{v+1}-1)(r(w)+1)-1$.
\end{proof}

\subsection{Tree Classes and Signed Systems}\label{app:cx-facts}

For the first proof we use the terminology of \cref{app:lf}: for a tree $t$, a
word $w$ \emph{meets} a leaf $v$ of $t$ if $v\pre w$; by the chain property
(\cref{sec:lf}), a word meets at most one leaf, a word that meets no leaf is a
position labeled $f$, and a proper prefix of a position is never a leaf; and
$w\in L^s_t$ iff $w$ meets a leaf of $t$ of sign $s$ (\cref{eq:lf-direct}).

\begin{proof}[Proof of \cref{lem:lf-classes}]
A word $w$ is a position labeled $f$ iff it meets no leaf, that is, iff
$w\notin L^0_t\cup L^1_t$. The positions of $t$ are $\eps$ and the children
of its positions labeled $f$, so $t$ has at most $1+nk$ positions if it has
$k$ positions labeled $f$; hence $t$ is finite iff
$\Sigma^*\setminus(L^0_t\cup L^1_t)$ is finite.

For regularity, let $w$ be a position of $t$ and $w^{-1}t$ the subtree of $t$
at $w$. The proper prefixes of $w$ are not leaves, so $wv$ meets a leaf of $t$
iff it meets a leaf $wv'$ with $v'$ a leaf of $w^{-1}t$, and the sign of
$wv'$ in $t$ is the sign of $v'$ in $w^{-1}t$ plus $\parity(w)$. Hence
\begin{equation}\label{eq:app-lf-quotient}
 L^s_{w^{-1}t}=w^{-1}L^{s\oplus\parity(w)}_t\qquad
 \text{for every position $w$ of $t$ and every sign $s$.}
\end{equation}
If $w$ is not a position of $t$, then a proper prefix of $w$ is a leaf, of
some sign $s$, and $w^{-1}L^s_t=\Sigma^*$ and $w^{-1}L^{1-s}_t=\varnothing$.
Suppose $L^0_t$ and $L^1_t$ are regular. By the Myhill--Nerode theorem they
have finitely many left quotients, so by \cref{eq:app-lf-quotient} the
subtrees of $t$ have finitely many leaf pairs, and by \cref{thm:lf-representation}(1) $t$ has finitely
many distinct subtrees. Conversely, if $t$ has finitely many distinct
subtrees, then by \cref{eq:app-lf-quotient} and the case of non-positions,
the pair $(w^{-1}L^0_t,w^{-1}L^1_t)$ takes finitely many values as $w$
ranges over all words, so both languages have finitely many left quotients
and are regular.
\end{proof}

\begin{proof}[Proof of \cref{lem:bg-signed-sat}]
If $\rho$ is a $\mathcal D$-solution of $\chi$, then $\rho^\pm$ is a
$\mathcal D$-solution of $\chi^\pm$ by \cref{lem:bg-signed}. For the converse
let $\sigma$ be a $\mathcal D$-solution of $\chi^\pm$.

For $n=0$ there are no children, so the literals of $\chi^\pm$ of sign $+$ in
\cref{tab:bg-signed} are the literals of $\chi$ with every variable $v$ replaced
by $v^+$, and the order is the same chain for every variance. Hence
$v\mapsto\sigma(v^+)$ solves $\chi$, and the three classes coincide.

Let $n\ge1$, and let $\chi^\flat$ be obtained by \cref{lem:bg-sandwich}, with
a fresh variable $y=y_{uv}$ for each copy literal $u\le v$ and the literals
$f(a_i)_i\le y$ and $y\le f(b_i)_i$. We extend $\sigma$ to a $\mathcal
D$-solution of $(\chi^\flat)^\pm$. The four signed literals of the two new
literals are
\[
 f(a_i^{c(i)})_i\le y^+,\quad y^-\le f(a_i^{\bar c(i)})_i,\quad
 y^+\le f(b_i^{c(i)})_i,\quad f(b_i^{\bar c(i)})_i\le y^- .
\]
Put $\sigma(y^+)=f(\sigma(a_i^{c(i)}))_i$ and
$\sigma(y^-)=f(\sigma(b_i^{\bar c(i)}))_i$. The first and fourth literals hold
with equality. The other two hold iff
$\sigma(a_i^{c(i)})\le\sigma(b_i^{c(i)})$ and
$\sigma(b_i^{\bar c(i)})\le\sigma(a_i^{\bar c(i)})$ for every $i$. For $i\ge1$
we have $a_i=b_i=u$, so both sides are equal. For $i=0$ these are the two inequalities
$\sigma(u^+)\le\sigma(v^+)$ and $\sigma(v^-)\le\sigma(u^-)$, in either case of
$c(0)$, and they hold because $u^+\le v^+$ and $v^-\le u^-$ are the signed
literals of $u\le v$. The new values lie in $\mathcal D$.

So $(\chi^\flat)^\pm$ is $\mathcal D$-satisfiable, and in particular satisfiable.
By \cref{thm:bg-nsse-sat}(2), $\chi^\flat$ has a solution, and a finite one
if $\mathcal D=\Trees_{\mathrm{fin}}$; by \cref{thm:bg-nsse-sat}(3) it has a
regular one. In each case $\chi^\flat$ has a $\mathcal D$-solution for the
variance $c$, whose restriction to the variables of $\chi$ is a
$\mathcal D$-solution of $\chi$ by \cref{lem:bg-sandwich}.

Finally, if $\chi$ has a solution, then so has $\chi^\flat$ by
\cref{lem:bg-sandwich}, hence a regular one by \cref{thm:bg-nsse-sat}(3),
and its restriction is a regular solution of $\chi$.
\end{proof}
